% Teil II: der durchgerechnete Weg.

% Gedankenschritt: Fakultät mit 2^{n-1} vergleichen, dann die Reihe als
% Teil + geometrische Reihe abschätzen.

\subsection*{Gegeben und gesucht}
(a) $n!>2^{n-1}$ f\"ur $n>2$; (b) $\sum_{n=2}^{\infty}1/n!<1$; (c) $2<e<3$;
(d) Python-Code f\"ur $s_n=\sum_{k=0}^{n}1/k!$ ($n=0,\dots,20$) und $e$ auf
12 Nachkommastellen.

\subsection*{a) $n!>2^{n-1}$ f\"ur $n>2$}
Induktion. IA ($n=3$): $3!=6>4=2^{2}$. \checkmark
IS: Ist $n!>2^{n-1}$, so
\begin{equation}
  (n+1)!=(n+1)\,n!>(n+1)\,2^{n-1}>2\cdot2^{n-1}=2^{n}.
\end{equation}
F\"ur $n+1\ge4$ gilt $n+1>2$. Damit gilt die Ungleichung f\"ur alle
$n>2$.

\subsection*{b) $\sum_{n=2}^{\infty}1/n!<1$}
F\"ur $n\ge3$ ist $1/n!<1/2^{n-1}$ (aus a)). Daher:
\begin{align}
  \sum_{n=2}^{\infty}\frac1{n!}
  &=\frac1{2!}+\sum_{n=3}^{\infty}\frac1{n!}
  <\frac12+\sum_{n=3}^{\infty}\frac1{2^{n-1}} \\
  &=\frac12+\Bigl(\frac14+\frac18+\cdots\Bigr)
  =\frac12+\frac12=1 .
\end{align}
Wegen strikt $<$ f\"ur alle $n\ge3$ ist die Summe strikt $<1$.

\subsection*{c) $2<e<3$}
Untere Schranke:
\begin{equation}
  e=1+1+\frac1{2!}+\frac1{3!}+\dots>2 .
\end{equation}
Obere Schranke: aus $e=2+\sum_{n=2}^\infty\frac1{n!}$ und b):
\begin{equation}
  e=2+\sum_{n=2}^{\infty}\frac1{n!}<2+1=3 .
\end{equation}

\subsection*{d) Python}
\begin{verbatim}
import math
for n in range(21):
    sn = sum([1/math.factorial(k) for k in range(n+1)])
    print(n, sn)
print(round(sum([1/math.factorial(k) for k in range(21)]), 12))
\end{verbatim}
Ausgabe (letzte Zeile): $e\approx2.\!718281828459$, also auf 12
Nachkommastellen
\begin{equation}
  e\approx2.\!718281828459 .
\end{equation}

\begin{ergebnis}
\begin{equation*}
  n!>2^{n-1}\ (n>2),\qquad
  \sum_{n=2}^{\infty}\frac1{n!}<1,\qquad
  2<e<3,\qquad
  e\approx2.\!718281828459 .
\end{equation*}
\end{ergebnis}

\begin{kontrolle}
\textbf{Teilsummen:} Bricht man bei $n=5$ ab: $\sum_{n=2}^5 1/n! =
\frac12+\frac16+\frac1{24}+\frac1{120}=0.7083<1$; die Konvergenz ist
weiter klar.

\textbf{e-Wert plausibilisieren:} $s_{20}=2.718281828459$ liegt im
Intervall $(2,3)$ wie in c) gezeigt; bekannte Dezimalentwicklung
$2.71828\dots$ stimmt \"uberein.

\textbf{Grenzfall/Laufzeit der Reihe:} $1/n!$ f\"allt riesigsteun schnell
(Im Nenner $n!$ w\"achst superlinear), deshalb konvergieren schon die
ersten 20 Glieder auf 12 Nachkommastellen. Das ist konsistent damit,
dass $e$ als Reihe \emph{schnell} konvergiert --- im Gegensatz zur
langsamen $1/k^2$-Reihe.
\end{kontrolle}

\begin{merksatz}
$e$ ist eine absolut (sogar superlinear) konvergente Reihe; die
Absch\"atzung \"uber $n!>2^{n-1}$ liefert Konvergenz und die Schranken
$2<e<3$. In Python z\"ahlt \texttt{math.factorial} die Reihe.
\end{merksatz}
